\section*{Hoofdstuk \ref{ch:b3:inorganic-materials} --- Anorganische materialen en nanomaterialen}

\begin{solution}{exo:b3:inorganic-materials:1}
$x^2 \propto t$: drie keer dikker duurt negen keer langer, \qty{90}{h}.
\end{solution}

\begin{solution}{exo:b3:inorganic-materials:2}
Hydrolyse: $\ce{Si(OC2H5)4 + 4H2O -> Si(OH)4 + 4C2H5OH}$.\\
Condensatie: $\ce{2Si(OH)4 -> Si2O(OH)6 + H2O}$.\\
Totaal: $\ce{Si(OC2H5)4 + 2H2O -> SiO2 + 4C2H5OH}$.
\end{solution}

\begin{solution}{exo:b3:inorganic-materials:3}
Vormers: $\ce{SiO2}$, $\ce{B2O3}$. Wijzigers: $\ce{Na2O}$, $\ce{CaO}$. Elk oxide-ion van $\ce{Na2O}$ verbreekt één brug Si--O--Si
tot twee uiteinden Si--O$^-$, elk gepaard met een $\ce{Na+}$: het netwerk wordt doorgeknipt,
de smelt vloeit gemakkelijker en smelt bij een lagere temperatuur.
\end{solution}

\begin{solution}{exo:b3:inorganic-materials:4}
Twaalf vijfhoeken (zoals voor elk \omterm{def:b3:inorganic-materials:carbon}{fullereen}). $V = 70$, $E = 3V/2 = 105$ bindingen, $F = 2 - V + E = 37$
vlakken, dus $37 - 12 = 25$ zeshoeken.
\end{solution}

\begin{solution}{exo:b3:inorganic-materials:5}
$\sqrt2\,(r_B + r_O) = 1.4142 \times \qty{200.5}{pm} = \qty{283.5}{pm}$. $\ce{SrTiO3}$: $284/283.5 = 1.00$, de ideale kubische perovskiet.
$\ce{BaTiO3}$: $301/283.5 = 1.06$, titaan te klein voor zijn octaëder, uit het centrum: \omterm{def:b3:inorganic-materials:ferroelectric}{ferro-elektrisch}.
$\ce{CaTiO3}$: $274/283.5 = 0.97$, calcium te klein voor zijn holte: gekantelde octaëders,
lagere symmetrie.
\end{solution}

\begin{solution}{exo:b3:inorganic-materials:6}
Si/Al $= 106/86 = 1.23$. $M = 86 \times 23.0 + 86 \times 27.0 + 106 \times 28.1 + 384 \times 16.0 = \qty{13422.6}{g/mol}$; 43 $\ce{Ca^2+}$ per
formule: $43/\qty{13422.6}{g/mol} = \qty{3.20}{mmol/g}$.
\end{solution}

\begin{solution}{exo:b3:inorganic-materials:7}
Afstand tussen naaste buren $d = a/\sqrt2 = \qty{288.4}{pm}$. Een cluster van $n$ schillen is $(2n + 1)d$ breed:
$2n + 1 \approx 5/0.2884 = 17.3$, dus $n = 8$ (\qty{4.9}{nm}). $N(8) = 2057$ atomen, $10 \times 64 + 2 = 642$ aan het oppervlak:
31\,\%.
\end{solution}

\begin{solution}{exo:b3:inorganic-materials:8}
$h^2/(8\mu R^2)$ met $\mu = \qty{9.109e-32}{kg}$: \qty{0.94}{eV} bij \qty{2.0}{nm} en \qty{0.42}{eV} bij \qty{3.0}{nm}. Emissie bij
$1.74 + 0.94 = \qty{2.68}{eV}$, \qty{463}{nm} (blauw), en $1.74 + 0.42 = \qty{2.16}{eV}$, \qty{574}{nm} (geel): de kleinere dot
zendt het blauwere licht uit.
\end{solution}

\begin{solution}{exo:b3:inorganic-materials:9}
$(10, 10)$: armstoel, $n - m = 0$, metallisch. $(10, 0)$: zigzag, 10 geen veelvoud van 3, halfgeleidend.
$(9, 0)$: zigzag, metallisch. $(12, 8)$: geen van beide (chiraal), $n - m = 4$, halfgeleidend.
\end{solution}

\begin{solution}{exo:b3:inorganic-materials:10}
$M = 4 \times 65.4 + 13 \times 16.0 + 24 \times 12.0 + 12 \times 1.0 = \qty{769.6}{g/mol}$; $a^3 = (\qty{25.82e-8}{cm})^3 = \qty{1.721e-20}{cm^3}$;
$\rho = 8 \times 769.6/(\num{6.022e23} \times \num{1.721e-20}) = \qty{0.594}{g.cm^{-3}}$. Een veld is \qty{7140}{m^2}: één gram bevat
$3000/7140 = 0.42$ veld, en een theelepel van enkele grammen meer dan een heel veld.
\end{solution}

\begin{solution}{exo:b3:inorganic-materials:11}
$\dfrac{10n^2 + 2}{(10n^3 + 15n^2 + 11n + 3)/3} = \dfrac{30n^2 + 6}{10n^3 + 15n^2 + 11n + 3} \to \dfrac{30n^2}{10n^3} = \dfrac3n$. Voor $n = 8$: exact
$642/2057 = 0.31$, tegenover $3/8 = 0.375$; de benadering overschat kleine
clusters, waarvan de binnenste schillen niet verwaarloosbaar zijn.
\end{solution}

\begin{solution}{exo:b3:inorganic-materials:12}
$2E = 3V$, $2E = 5p + 6h + 7s$, $F = p + h + s$. Euler: $V - E + F = 2$, met $V = 2E/3$: $F - E/3 = 2$, dus
$p + h + s - (5p + 6h + 7s)/6 = 2$, d.w.z.\ $(p - s)/6 = 2$: $p - s = 12$. Elke zevenhoek vereist een extra
vijfhoek.
\end{solution}

\begin{solution}{pb:b3:inorganic-materials:1}
\textbf{1.} $\mathrm{Na_{12}Al_{12}Si_{12}O_{48}}$: $12 \times 23.0 + 12 \times 27.0 + 12 \times 28.1 + 48 \times 16.0 = \qty{1705.2}{g/mol}$.

\textbf{2.} $1705.2 + 27 \times 18.0 = \qty{2191.2}{g/mol}$.

\textbf{3.} Si/Al $= 12/12 = 1$.

\textbf{4.} Geen verbindingen Al--O--Al; met Si/Al $= 1$ heeft elk Si vier Al-buren en elk
Al vier Si: strikte afwisseling.

\textbf{5.} $-12$, één per aluminiumatoom.

\textbf{6.} $486.0/2191.2 = 22.2$\,\%.

\textbf{7.} $\ce{2Na+}(\text{zeoliet}) + \ce{Ca^2+}(\text{aq}) \rightleftharpoons \ce{Ca^2+}(\text{zeoliet}) + \ce{2Na+}(\text{aq})$.

\textbf{8.} Zes (twaalf ladingen).

\textbf{9.} $6/\qty{1705.2}{g/mol} = \qty{3.52}{mmol/g}$.

\textbf{10.} $6/\qty{2191.2}{g/mol} = \qty{2.74}{mmol/g}$.

\textbf{11.} $\qty{3.52}{mmol/g} \times \qty{100.1}{g/mol} = \qty{352}{mg}$ $\ce{CaCO3}$ per gram.

\textbf{12.} $\qty{2.0}{mmol/L} \times \qty{15}{L} = \qty{30}{mmol}$.

\textbf{13.} $\qty{30}{mmol}/\qty{2.74}{mmol/g} = \qty{11}{g}$.

\textbf{14.} $\qty{10.9}{g}/0.60 = \qty{18}{g}$.

\textbf{15.} $\qty{60}{mmol}$ $\ce{Na+}$, $\qty{0.060}{mol} \times \qty{23.0}{g/mol} = \qty{1.4}{g}$.

\textbf{16.} Fosfaten in afvalwater voeden algen in rivieren en meren; de
algenbloei en de afbraak ervan verbruiken de opgeloste zuurstof
(eutrofiëring).

\textbf{17.} Het kleinere $\ce{Mg^2+}$ houdt zijn watermoleculen steviger vast; het moet ze afstoten
om door het venster te gaan en de plaatsen te bereiken, wat bij
wastemperaturen traag verloopt.

\textbf{18.} Het naakte ion is ongeveer \qty{0.20}{nm} breed, de helft van het venster.

\textbf{19.} Het gehydrateerde ion is groter dan het venster: het moet een
deel van zijn watermantel afstoten om erdoor te gaan, en de
zuurstofatomen van het raamwerk nemen de plaats van het water in.

\textbf{20.} Het raamwerk is negatief geladen en wisselt alleen kationen
uit; een groot anion met een lange keten wordt afgestoten en kan niet
door het venster, zodat het in het water blijft, waar het zijn werk doet.

\textbf{21.} Gedehydrateerd neemt de \omterm{def:b3:inorganic-materials:zeolite}{zeoliet} water sterk op in zijn kooien;
alleen moleculen die kleiner zijn dan het venster komen binnen, zodat hij
moleculen volgens grootte scheidt: een zeef op moleculaire schaal.

\textbf{22.} De grotere ionen $\ce{K+}$ zitten in de vensters en vernauwen ze: alleen kleinere
moleculen, zoals water, worden toegelaten.

\textbf{23.} Het slanke, lineaire isomeer \textit{para} past in de kanalen en diffundeert snel; de
bredere isomeren \textit{ortho} en \textit{meta} diffunderen traag, blijven achter en isomeriseren.

\textbf{24.} Watervrij \omterm{def:b3:inorganic-materials:zeolite}{zeoliet} A kan in theorie \qty{3.52}{mmol} $\ce{Ca^2+}$ per gram opnemen.
\end{solution}
